\documentclass[11pt]{article}
\title{Nonconsecutive Subset Sums Modulo Three: Exact Encodings and an Unproved Recurrence Candidate}
\date{}
\author{SIGMA AI-assisted research synthesis}
\usepackage[T1]{fontenc}
\usepackage[utf8]{inputenc}
\usepackage[margin=1in]{geometry}
\usepackage{amsmath,amssymb,booktabs,longtable}
\usepackage{microtype}
\usepackage{titlesec}
\emergencystretch=2em
\raggedbottom
\linespread{1.07}
\usepackage{xcolor}
\definecolor{blue}{HTML}{24485E}
\titleformat{\section}{\Large\bfseries\sffamily\color{blue}}{\thesection}{0.8em}{}
\titleformat{\subsection}{\large\bfseries\sffamily}{\thesubsection}{0.7em}{}
\setlength{\parskip}{.6em}
\setlength{\parindent}{0pt}
\begin{document}
\maketitle

\begin{abstract}
How many subsets of an initial interval contain no consecutive elements and have sum divisible by three? We examine this question through weighted Fibonacci polynomials, residue-state counting, and a reversible gap-removal transformation. The weighted formulation separates the desired count into an ordinary Fibonacci contribution and two nontrivial cube-root contributions. Exact finite calculations give $a_{12}=128$ and disprove the tempting rule $a_{n+3}=a_{n+2}+a_n$. A different integer-coefficient recurrence of order twelve agrees with the reported finite tests, but its phase-ordered matrix derivation and boundary cases remain incomplete.

The gap-removal bijection provides a second encoding, with a corrected Gaussian-binomial normalization explaining the weight $q^{k^2}$. These results isolate the remaining mathematical tasks: a global recurrence proof, sufficient justified initial values, and a growth bound for the nontrivial cube-root terms. In particular, the proposed limiting density of one third remains unproved by the available argument.
\end{abstract}
\section*{Evidence coverage limitation}
The evidence index omitted 996 material saved research field fragment(s) across 6 source(s). These fragments were not reviewed for this paper. Interpretations that could depend on omitted material remain provisional; the complete frozen source snapshot and omission ledger are retained with the generation record.
\section{Introduction}
For each integer $n\geq 0$, let $a_n$ count subsets of $\{1,\ldots,n\}$ containing no consecutive integers and having sum congruent to zero modulo three. Indices are one-based, and the empty subset is included, so $a_0=1$. The target is a rigorously justified constant-coefficient recurrence with integer coefficients, sufficient initial values, the exact value of $a_{12}$, and the limit $a_n/F_{n+2}\to 1/3$, where $F_0=0$ and $F_1=1$.

The weighted enumeration connects the local adjacency restriction to the global congruence condition. Its specialization at one counts all admissible subsets, while evaluations at primitive cube roots isolate the residue-zero count. This separation explains why a recurrence proof and an asymptotic-density proof require different arguments: an annihilating relation alone does not establish that the nontrivial contributions are negligible.

\section{Evidence and scope}
Finite counting used states consisting of the current sum residue and a bit recording whether the preceding integer was selected. Independent bit-mask enumeration supplied the small cases, while exact-integer dynamic programming extended the count through $n=60$ and checked the proposed recurrence on the corresponding finite range. These calculations served as exact finite evaluations and counterexample searches, not as a proof of a relation for arbitrary indices.

The structural analysis compared the weighted Fibonacci recurrence with a gap-removal bijection for subsets of fixed cardinality. Exact integer-polynomial comparisons through $n=18$ corroborated the Gaussian-binomial expression. A normalization correction distinguished the positional shift from the baseline weight in the standard Gaussian polynomial; the transfer-matrix argument and the asymptotic estimate were assessed separately because their global proof obligations remained open.

\section{Weighted enumeration and exact finite counts}
Weighted enumeration separates the adjacency restriction from the congruence condition and gives an exact route to finite counts without assuming a constant-coefficient recurrence. For $n\geq 0$, let $\mathcal A_n$ be the family of subsets of $\{1,\ldots,n\}$ containing no two consecutive integers, and define $P_n(q)=\sum_{S\in\mathcal A_n}q^{\sum_{i\in S}i}$. The desired number is $a_n=\#\{S\in\mathcal A_n:\sum_{i\in S}i\equiv0\pmod 3\}$. The polynomial records the full integer sum of each admissible subset, retaining the information needed to extract its residue class.

The boundary polynomials are $P_0(q)=1$ and $P_1(q)=1+q$. For the empty interval, the only subset is empty, its sum is $0$, and its contribution is therefore $1$. For the one-element interval, the empty subset and $\{1\}$ contribute $1$ and $q$, respectively. These boundary values specify the enumeration directly; they do not require extending a recurrence to negative indices.

For $n\geq2$, partition $\mathcal A_n$ according to whether its members contain $n$. Subsets omitting $n$ are precisely the members of $\mathcal A_{n-1}$ and contribute $P_{n-1}(q)$. A subset containing $n$ must omit $n-1$; removing $n$ gives a member of $\mathcal A_{n-2}$, and adjoining $n$ reverses this operation while increasing the sum by $n$. Thus the disjoint partition proves $P_n(q)=P_{n-1}(q)+q^nP_{n-2}(q)$ for $n\geq2$.

Specialization at $q=1$ discards the sum weights and counts all admissible subsets. The resulting recurrence is $P_n(1)=P_{n-1}(1)+P_{n-2}(1)$, with $P_0(1)=1$ and $P_1(1)=2$. With Fibonacci numbers defined by $F_0=0$, $F_1=1$, and $F_{m+2}=F_{m+1}+F_m$ for $m\geq0$, these initial values and the recurrence give $P_n(1)=F_{n+2}$. Consequently, the proposed denominator is exactly the total number of admissible subsets, rather than an approximation to that number.

To isolate sums divisible by three, let $\omega$ be a primitive cube root of unity, so that $\omega^3=1$ and $1+\omega+\omega^2=0$. For an integer sum $s$, the average $(1+\omega^s+\omega^{2s})/3$ equals $1$ if $s\equiv0\pmod3$ and $0$ otherwise. Applying this identity to each monomial in $P_n$ gives $a_n=(P_n(1)+P_n(\omega)+P_n(\omega^2))/3$. Evaluation at $\omega$ alone is not coefficient extraction: it assigns different complex weights to the three residue classes, whereas the average cancels the two unwanted classes.

A residue-and-flag dynamic program gives a second exact encoding that uses only integer counts. Let $D_i(r,b)$ count admissible subsets of $\{1,\ldots,i\}$ whose sum has residue $r\in\{0,1,2\}$, where $b\in\{0,1\}$ records whether $i$ is selected. Initialize $D_0(0,0)=1$ and all other states to $0$, with the flag at the empty interval interpreted as allowing the first integer to be selected. This initialization counts the empty subset once and places it in residue class $0$.

When processing $i\geq1$, omission of $i$ is allowed from either previous flag, so $D_i(r,0)=D_{i-1}(r,0)+D_{i-1}(r,1)$. Inclusion is allowed only when $i-1$ was omitted and increases the sum by $i$, giving $D_i(r,1)=D_{i-1}(r-i,0)$, where residue arguments are reduced modulo $3$. These transitions are exhaustive and disjoint: every admissible subset either omits $i$, or contains it and has an admissible predecessor with the preceding index unselected. The required count is therefore $a_i=D_i(0,0)+D_i(0,1)$.

The reported exact integer computation with these six states gives $a_{12}=128$. Its force is finite and precise: the value is obtained by counting admissible subsets through exhaustive state transitions, not by extrapolating a recurrence fitted to earlier terms. Together, the weighted polynomial and the residue-and-flag encoding establish the Fibonacci normalization, an exact congruence filter, and the finite count independently of the proposed higher-order recurrence.

\section{Gap removal and the Gaussian-binomial weight}
Removing the mandatory gaps between selected positions gives an exact second encoding of the weighted enumeration and explains the exponent $k^2$ in its Gaussian-binomial form. The essential distinction is between two offsets: the gaps removed from an admissible subset and the minimum sum of an ordinary subset. Keeping these offsets separate corrects the normalization without changing the resulting formula.

Fix $n\geq 0$ and an admissible subset of positive cardinality $k$, written as $1\leq i_1<\cdots<i_k\leq n$, with $i_{j+1}-i_j\geq 2$. Set $y_j=i_j-(j-1)$. Then $y_{j+1}-y_j=i_{j+1}-i_j-1\geq 1$, while $y_1=i_1\geq 1$ and $y_k=i_k-(k-1)\leq n-k+1$. Thus the transformed positions form an ordinary $k$-element subset of $\{1,\ldots,n-k+1\}$. The original gap condition also gives $i_k\geq 2k-1$, so the possible cardinalities satisfy $k\leq\lceil n/2\rceil$.

Conversely, start with any ordinary subset $1\leq y_1<\cdots<y_k\leq n-k+1$ and define $i_j=y_j+(j-1)$. Its endpoints satisfy $i_1\geq 1$ and $i_k\leq n$, and its successive gaps satisfy $i_{j+1}-i_j=y_{j+1}-y_j+1\geq 2$. It therefore produces an admissible subset of the original interval. The two transformations undo one another coordinate by coordinate, establishing a bijection for every allowable positive cardinality; at cardinality $k=0$, the empty subset corresponds to itself.

The gap-removal map changes the sum of the selected positions by a fixed amount depending only on cardinality: $\sum_{j=1}^{k}i_j=\sum_{j=1}^{k}y_j+\sum_{j=1}^{k}(j-1)=\sum_{j=1}^{k}y_j+k(k-1)/2$. Consequently, the weight enumerator for admissible subsets of cardinality $k$ is $q^{k(k-1)/2}$ times the ordinary subset enumerator weighted by $q^{\sum_j y_j}$. This factor accounts only for the removed gaps; it does not yet include the normalization of the Gaussian polynomial.

For integers $0\leq k\leq m$, use the partition-normalized Gaussian polynomial convention \$\textbackslash{}displaystyle \{m\textbackslash{}brack k\}\_q=\textbackslash{}sum\_\{1\textbackslash{}leq y\_1\textless{}\textbackslash{}cdots\textless{}y\_k\textbackslash{}leq m\}q\textasciicircum{}\{\textbackslash{}sum\_\{j=1\}\textasciicircum{}\{k\}y\_j-k(k+1)/2\}\$, with \$\{m\textbackslash{}brack 0\}\_q=1\$. The baseline $k(k+1)/2$ is the sum of the smallest ordinary subset, whose positions are $1,\ldots,k$. Hence the unnormalized ordinary subset enumerator is \$\textbackslash{}displaystyle \textbackslash{}sum\_\{1\textbackslash{}leq y\_1\textless{}\textbackslash{}cdots\textless{}y\_k\textbackslash{}leq m\}q\textasciicircum{}\{\textbackslash{}sum\_\{j=1\}\textasciicircum{}\{k\}y\_j\}=q\textasciicircum{}\{k(k+1)/2\}\{m\textbackslash{}brack k\}\_q\$. In particular, its prefactor is $q^{k(k+1)/2}$, not $q^{k(k-1)/2}$.

Taking $m=n-k+1$ and combining the two distinct shifts gives $k(k-1)/2+k(k+1)/2=k^2$. Summing the fixed-cardinality contributions therefore yields the polynomial identity \$\textbackslash{}displaystyle P\_n(q)=\textbackslash{}sum\_\{k=0\}\textasciicircum{}\{\textbackslash{}lceil n/2\textbackslash{}rceil\}q\textasciicircum{}\{k\textasciicircum{}2\}\{n-k+1\textbackslash{}brack k\}\_q\$ for every $n\geq 0$. The term with $k=0$ is exactly $1$, representing the empty subset; in particular, the formula also gives $P_0(q)=1$. Because the bijection preserves the weights with the stated offset, this is an identity of integer-coefficient polynomials before any specialization of $q$.

Reported exact coefficient comparisons through $n=18$ agree with this expression, but their role is corroborative: the reversible map and the explicit normalization establish the identity for arbitrary $n$. The correction concerns the explanation of the exponent, not the final factor $q^{k^2}$. This supplies the central argument with an exact cardinality-resolved encoding, while leaving the passage from that encoding to a constant-coefficient recurrence as a separate algebraic task.

\section{A refuted short rule and the order-twelve candidate}
The exact counterexample to the short recurrence separates a disproved pattern from the still viable order-twelve candidate. The small counts, including the empty subset, are $(a_0,a_1,\ldots,a_8)=(1,1,1,2,3,5,8,12,19)$. Although $a_{n+3}=a_{n+2}+a_n$ holds at $n=0$ and $n=1$, it fails at $n=2$: $a_5=5$, whereas $a_4+a_2=3+1=4$. This is a refutation of that proposed rule, not merely an absent proof, and it establishes no replacement recurrence.

The distinct candidate is $a_{n+12}-5a_{n+9}+2a_{n+6}+5a_{n+3}+a_n=0$, proposed for every $n\geq 0$. An exact six-state dynamic program, recording the sum modulo three and whether the preceding integer was selected, computed counts through $n=60$ and tested this relation at every starting index from $n=0$ through $n=48$. All reported residuals vanished. These bounded integer calculations corroborate the candidate, but extending the range of agreement does not turn it into a consequence of the counting definition.

The proposed algebraic explanation uses the decomposition $3a_n=P_n(1)+P_n(\omega)+P_n(\omega^2)$, where $\omega$ is a primitive cube root of unity and $P_n(1)=F_{n+2}$. Its essential step is to advance each component by three indices, rather than to guess a recurrence directly for their sum. The proposed three-step block relation for each nontrivial evaluation has characteristic polynomial $x^2-x-1$, while the stated relation for the ordinary Fibonacci component has characteristic polynomial $x^2-4x-1$. Thus the relevant assertion concerns the three residue-class subsequences of each component, not just a few successive values of $a_n$.

To make the conditional combination precise, let $T$ denote the three-step shift, so that $(Tu)_n=u_{n+3}$. If $T^2-T-1$ annihilates both nontrivial components and $T^2-4T-1$ annihilates the Fibonacci component, their product annihilates the sum: each component is killed by one factor, and the factors commute because both are polynomials in the same shift. The polynomial identity $(x^2-x-1)(x^2-4x-1)=x^4-5x^3+2x^2+5x+1$ then gives exactly the displayed candidate. No additional copy of the first factor is needed merely because there are two nontrivial evaluations; the proposed argument requires that the same factor annihilate each of them.

The unresolved step is the derivation of those component relations from the weighted transition, not the multiplication of the two polynomials. At a cube root of unity, the coefficients in $P_n(q)=P_{n-1}(q)+q^nP_{n-2}(q)$ are periodic with period three. The proposed block calculation asserts trace $1$ and determinant $-1$ for the nontrivial three-step transfer, which would supply the polynomial $x^2-x-1$. However, the exact ordered products have not been established for all three starting phases. Because the individual transfer matrices do not generally commute, periodicity alone does not justify replacing the phase-dependent products by an unstated common matrix identity.

The starting boundary is a separate part of this same proof obligation. The weighted recurrence is stated for $n\geq 2$, with the initial polynomials at $n=0$ and $n=1$ handled separately; consequently, a block argument cannot simply apply that recurrence across the initial boundary without checking its validity there. A global derivation must establish the required identity in every phase and verify any exceptional initial indices directly. This matters because the candidate is asserted beginning at $n=0$, rather than only after an unspecified transient range.

A complete recurrence specification would also require sufficient justified initial values. For the displayed order-twelve relation beginning at $n=0$, the values $a_0,\ldots,a_{11}$ would determine all subsequent terms; equivalently, each of the three residue-class subsequences would require four starting values. The explicit small-value list above reaches only $a_8$, so it is not by itself such a specification. The evidence therefore replaces the false short rule with a precise algebraic target: an ordered three-step identity whose component polynomials explain the candidate’s coefficients, rather than a longer fit whose finite agreement substitutes for that identity.

\section{The remaining asymptotic argument}
The asymptotic question reduces to controlling the discrepancy contributed by the two nontrivial cube-root evaluations, rather than estimating the total number of admissible subsets anew. Let $\omega$ be a primitive cube root of unity and retain the weighted enumeration polynomial $P_n(q)$. The root-of-unity decomposition and the identity $P_n(1)=F_{n+2}$ give the exact relation $\displaystyle \frac{a_n}{F_{n+2}}-\frac13=\frac{P_n(\omega)+P_n(\omega^2)}{3F_{n+2}}$. Thus the Fibonacci contribution already supplies the proposed main term; the remaining issue is whether the weighted contributions become negligible after normalization.

The two nontrivial evaluations are complex conjugates, so the discrepancy can also be written as $\displaystyle \frac{2\operatorname{Re}P_n(\omega)}{3F_{n+2}}$. In particular, proving $P_n(\omega)=o(F_{n+2})$ would control both terms and establish the requested limit. This is a sufficient bound on the full complex evaluation; the exact error identity requires only its real part to be negligible. Neither the exact finite count nor agreement with a proposed recurrence supplies this growth estimate.

The proposed route is to examine the weighted transfer relation over blocks of three indices. For each fixed residue $r\in\{0,1,2\}$, the relevant quantities are the subsequences $P_{3m+r}(\omega)$ and $P_{3m+r}(\omega^2)$, whose growth must be compared with $F_{3m+r+2}$. A spectral argument would first identify the correctly ordered block matrix for each phase and then use its eigenvalues to bound the powers acting on the corresponding initial state. Because the phase matrices do not generally commute, periodicity alone does not justify replacing these ordered products by an unspecified common block.

The suggested algebraic mechanism assigns trace $1$ and determinant $-1$ to the nontrivial three-step blocks, leading to the proposed characteristic polynomial $x^2-x-1$; the corresponding polynomial for the ordinary Fibonacci component is $x^2-4x-1$. These statements describe the intended bridge to a spectral comparison, not an established estimate for the weighted subsequences. The ordered product identities still require justification for every phase, with the initial boundary handled separately because the weighted recurrence is stated only for $n\geq2$.

Once that bridge is proved, the remaining asymptotic task is to show that the nontrivial modes grow strictly more slowly than the Fibonacci component and that the resulting bounds apply to the actual initial states.

Control on all three residue subsequences would then yield control of the full sequence, since every index has the form $n=3m+r$ with $r\in\{0,1,2\}$. This explains why a global recurrence proof and a density proof are related but distinct goals: a recurrence constrains possible modes, whereas the limit also requires their growth and their contributions to the specified solution to be controlled. Finite recurrence residuals or observed ratios cannot replace that comparison.

The achieved results therefore have different logical roles. The reported exact value $a_{12}=128$ answers the finite counting question, while the weighted formulation and reversible gap-removal transformation give structural encodings; the order-twelve relation remains a candidate rather than a globally justified recurrence with sufficient established initial data. For the density question, the root-of-unity decomposition makes the outstanding requirement precise: the proposed one-third main term becomes a limiting theorem only when the nontrivial weighted contributions are shown to be negligible relative to $F_{n+2}$.

\section{Verification status and limitations}
\begin{itemize}
\item The recurrence tests and polynomial comparisons cover finite ranges only. Their role is corroboration or refutation, not establishment of universal identities.
\item The evidence contains an inconsistency in proof status for the Gaussian-binomial identity: one assessment accepts the combinatorial identity, while later objections still request a complete convention-specific symbolic derivation. The exact bijection and corrected normalization are retained without treating that disagreement as resolved.
\item The computational and argumentative reports have not been independently reproduced or formally verified. Their fixture provenance and assessment agreement do not provide independent validation.
\item The supplied index is a bounded selection rather than an exhaustive review of all saved notes. Absence of a complete initial-value table or proof in this preview is not evidence that no further relevant material exists.
\end{itemize}
\section{Open obligations}
\begin{itemize}
\item Complete the all-index Gaussian-binomial derivation under the stated partition-normalized convention, including the cardinality endpoints and the empty-subset term.
\item Derive the ordered three-step transfer products for every phase and verify the proposed trace, determinant, and component polynomial identities symbolically.
\item Show precisely how the component annihilating relations imply the candidate scalar recurrence for every claimed index, with a separate treatment of the boundary at $n=0$.
\item Supply and justify sufficient initial values for the selected recurrence. The indexed small-value list alone does not initialize the displayed order-twelve relation.
\item Prove a growth bound making the two nontrivial cube-root contributions negligible relative to $F_{n+2}$, and then deduce the density limit from the filtering identity.
\end{itemize}
\section{Conclusion}
The weighted recurrence, residue-state representation, and gap-removal bijection provide complementary exact descriptions of the counting problem. The corrected normalization explains why the proposed Gaussian-binomial expression carries the factor $q^{k^2}$.

The reported exact finite value is $a_{12}=128$. The short rule $a_{n+3}=a_{n+2}+a_n$ is definitively refuted by the available small counts.

The order-twelve recurrence remains a supported candidate rather than a proved theorem. Its global transfer-matrix derivation, sufficient initialization, and the asymptotic-density argument remain separate obligations; the available evidence therefore constitutes partial progress rather than a complete solution.

\section*{Claim status and evidence}
\subsection{C176 -- Established}
The counting problem uses one-based indices, excludes consecutive selected integers, includes the empty subset, and asks for an integer-coefficient constant-coefficient recurrence, sufficient initial values, $a_{12}$, and the limit $a_n/F_{n+2}\to 1/3$.
These are the problem's definitions and requested conclusions, not assertions that the conclusions have been proved. The empty subset fixes the boundary value $a_0=1$.
\textit{Evidence: }project:57d67c57-15dc-433d-828a-7953f7986ff8, artifact:1a9a0c44-48ae-49e7-8bea-0e8d47297f9d:1cc6ec50-87c0-46f4-b1dd-a82857d74f9e
\subsection{C177 -- Established}
For $P_n(q)=\sum_S q^{\sum_{i\in S}i}$, where the sum runs over nonconsecutive subsets of $\{1,\ldots,n\}$, the weighted recurrence is $P_n(q)=P_{n-1}(q)+q^nP_{n-2}(q)$ for $n\geq2$, with $P_0(q)=1$ and $P_1(q)=1+q$. Moreover, $P_n(1)=F_{n+2}$ and $a_n=(P_n(1)+P_n(\omega)+P_n(\omega^2))/3$ for a primitive cube root $\omega$.
Subsets omitting the largest integer contribute the first term. Subsets containing it must omit its predecessor, leaving an admissible subset of the shorter interval and contributing the additional weight $q^n$. At one, this is the Fibonacci counting recurrence with the appropriate two initial values. Averaging the three root-of-unity evaluations retains precisely those monomials whose exponents are divisible by three.
\textit{Evidence: }strategy:00469f7f-ba53-4a12-944d-2d3295e48d1f, artifact:1aabbfc0-e642-4581-a96f-2e33ac17ee19:ccc1045f-479d-4018-b8ed-42a60b94e567, session:1aabbfc0-e642-4581-a96f-2e33ac17ee19
\subsection{C178 -- Established}
An exact six-state counting representation records a residue $r\in\{0,1,2\}$ and a previous-selection bit $b\in\{0,1\}$. Omitting the next integer resets the bit, while including it is permitted only from a state with $b=0$ and updates the residue by that integer modulo three.
The bit retains exactly the information needed to enforce adjacency, and the residue retains exactly the information needed for the congruence condition. The transitions depend on the processed integer only modulo three, giving a periodic transfer description without yet proving the proposed scalar recurrence.
\textit{Evidence: }strategy:5c8ce76f-4965-434b-8862-841c92de89e4, artifact:6c552b3c-44e4-4b57-8b93-e5a645d057d1:5e15b337-adf1-483c-8ba6-7cd2de5f500b
\subsection{C179 -- Established}
The rule $a_{n+3}=a_{n+2}+a_n$ is false. The reported exact values are $(a_0,\ldots,a_8)=(1,1,1,2,3,5,8,12,19)$, and at $n=2$ they give $a_5=5\neq a_4+a_2=4$.
A single exact counterexample refutes a universal recurrence. Agreement at the first two transitions does not survive the next transition and therefore provides no basis for extending this short rule.
\textit{Evidence: }artifact:1a9a0c44-48ae-49e7-8bea-0e8d47297f9d:1cc6ec50-87c0-46f4-b1dd-a82857d74f9e, session:1a9a0c44-48ae-49e7-8bea-0e8d47297f9d
\subsection{C180 -- Established}
For a nonconsecutive subset $1\leq i_1<\cdots<i_k\leq n$, the transformation $y_j=i_j-(j-1)$ is a bijection to ordinary $k$-subsets of $\{1,\ldots,n-k+1\}$, and $\sum_j i_j=\sum_j y_j+k(k-1)/2$.
Successive original positions differ by at least two, so subtracting successive offsets makes the transformed positions strictly increasing. The stated endpoint bounds follow from the first and last positions. The inverse transformation $i_j=y_j+(j-1)$ restores gaps of at least two and the original interval bounds, while summing the offsets gives the weight shift.
\textit{Evidence: }strategy:9ad51640-ff9a-4d22-a928-518f032003bf, artifact:f7301d22-fd90-4675-95cd-17b006262ef9:02ca743b-d66f-41c0-8692-34579ea5adfa, artifact:7c70ca6e-3fb2-4cda-b75a-72702ee10d97:516c2288-3458-40b2-9e02-87a918c35a3a
\subsection{C181 -- Established}
Under the standard partition-normalized Gaussian-binomial convention, the ordinary $k$-subset enumerator has baseline weight $q^{k(k+1)/2}$. Combining it with the gap-removal shift $q^{k(k-1)/2}$ gives $q^{k^2}$; using the latter shift for both factors is an explanatory error.
The normalized Gaussian polynomial removes the minimal sum $1+\cdots+k=k(k+1)/2$ from ordinary subset positions. Restoring that baseline and then restoring the removed gaps adds the two distinct offsets. Their sum is $k^2$, preserving the displayed formula while correcting its earlier explanation.
\textit{Evidence: }artifact:7c70ca6e-3fb2-4cda-b75a-72702ee10d97:516c2288-3458-40b2-9e02-87a918c35a3a, session:7c70ca6e-3fb2-4cda-b75a-72702ee10d97, artifact:f7301d22-fd90-4675-95cd-17b006262ef9:02ca743b-d66f-41c0-8692-34579ea5adfa
\subsection{C182 -- Unresolved}
The proposed identity is $P_n(q)=\sum_{k=0}^{\lceil n/2\rceil}q^{k^2}\genfrac{[}{]}{0pt}{}{n-k+1}{k}_q$, with the term for $k=0$ equal to one. The evidence supplies the bijective mechanism and corrected normalization, but retains an explicit obligation to complete the convention-specific symbolic derivation for all indices.
Some summaries describe the formula as a correct combinatorial identity, whereas later objections still request its complete symbolic proof under the specified convention. The plan preserves the supported bijection and normalization separately rather than treating finite coefficient agreement or reviewer agreement as completed verification of the full presentation.
\textit{Evidence: }artifact:f7301d22-fd90-4675-95cd-17b006262ef9:02ca743b-d66f-41c0-8692-34579ea5adfa, artifact:7c70ca6e-3fb2-4cda-b75a-72702ee10d97:516c2288-3458-40b2-9e02-87a918c35a3a
\subsection{C183 -- Established}
The reported exact finite calculations give $a_{12}=128$.
Both the residue-counting and weighted-enumeration records report this finite value. It is a finite counting result, not a consequence of an already proved order-twelve recurrence. Its status does not imply independent reproduction or formal verification of the underlying computation.
\textit{Evidence: }artifact:1aabbfc0-e642-4581-a96f-2e33ac17ee19:ccc1045f-479d-4018-b8ed-42a60b94e567, artifact:f7301d22-fd90-4675-95cd-17b006262ef9:02ca743b-d66f-41c0-8692-34579ea5adfa, session:1aabbfc0-e642-4581-a96f-2e33ac17ee19
\subsection{C184 -- Conjecture}
The candidate integer-coefficient recurrence is $a_{n+12}-5a_{n+9}+2a_{n+6}+5a_{n+3}+a_n=0$ for $n\geq0$.
The proposed three-step argument assigns the nontrivial cube-root components the polynomial $x^2-x-1$ and the ordinary Fibonacci component the polynomial $x^2-4x-1$. Their product has the coefficients displayed in the candidate recurrence. The common annihilation argument is conditional because the ordered matrix identities and boundary coverage have not been completed.
\textit{Evidence: }artifact:1aabbfc0-e642-4581-a96f-2e33ac17ee19:ccc1045f-479d-4018-b8ed-42a60b94e567, artifact:6c552b3c-44e4-4b57-8b93-e5a645d057d1:5e15b337-adf1-483c-8ba6-7cd2de5f500b
\subsection{C185 -- Established}
The reported checks include exact integer-polynomial comparisons through $n=18$, exact residue-state counts through $n=60$, and vanishing candidate-recurrence residuals through $n=48$. These are bounded corroboration, not proofs for arbitrary indices.
The bounds and exact-arithmetic descriptions are explicitly reported in the indexed evidence. Finite agreement can reveal a counterexample within its range but cannot establish the proposed identities outside that range.
\textit{Evidence: }artifact:f7301d22-fd90-4675-95cd-17b006262ef9:02ca743b-d66f-41c0-8692-34579ea5adfa, artifact:7c70ca6e-3fb2-4cda-b75a-72702ee10d97:516c2288-3458-40b2-9e02-87a918c35a3a, artifact:6c552b3c-44e4-4b57-8b93-e5a645d057d1:5e15b337-adf1-483c-8ba6-7cd2de5f500b
\subsection{C186 -- Unresolved}
A global proof of the order-twelve recurrence and a sufficient justified initialization remain incomplete. The outstanding tasks include ordered three-step products for every phase, the proposed trace and determinant calculation, and explicit treatment of the boundary at $n=0$.
The transfer matrices need not commute, so periodicity alone does not justify an unspecified product order. The weighted recurrence is stated only for $n\geq2$, making its use at the earliest boundary a separate obligation. The indexed preview also does not supply a complete justified list of the twelve consecutive initial values needed for the displayed forward recurrence.
\textit{Evidence: }artifact:1aabbfc0-e642-4581-a96f-2e33ac17ee19:ccc1045f-479d-4018-b8ed-42a60b94e567, artifact:6c552b3c-44e4-4b57-8b93-e5a645d057d1:5e15b337-adf1-483c-8ba6-7cd2de5f500b, artifact:7c70ca6e-3fb2-4cda-b75a-72702ee10d97:516c2288-3458-40b2-9e02-87a918c35a3a
\subsection{C187 -- Unresolved}
The limit $a_n/F_{n+2}\to1/3$ remains unproved in the available argument. The identified missing estimate is that the nontrivial cube-root contributions are negligible relative to $F_{n+2}$.
The root-of-unity filter expresses the difference from one third using the two nontrivial evaluations. An analytic growth bound would therefore connect the weighted decomposition to the requested limit. Neither finite ratios nor agreement with the proposed recurrence supplies that bound.
\textit{Evidence: }artifact:1aabbfc0-e642-4581-a96f-2e33ac17ee19:ccc1045f-479d-4018-b8ed-42a60b94e567, artifact:f7301d22-fd90-4675-95cd-17b006262ef9:02ca743b-d66f-41c0-8692-34579ea5adfa, artifact:6c552b3c-44e4-4b57-8b93-e5a645d057d1:5e15b337-adf1-483c-8ba6-7cd2de5f500b
\subsection{C188 -- Limitation}
The indexed mathematical arguments and computational reports are not independently verified; several are explicitly deterministic fixture narratives, and agreement between their associated assessments is not independent validation.
The sources explicitly distinguish fixture narratives and finite calculations from completed mathematical verification. Established labels in this register identify definitions, supported structural steps, counterexamples, or reported finite results; they do not certify the full requested theorem.
\textit{Evidence: }session:1a9a0c44-48ae-49e7-8bea-0e8d47297f9d, session:1aabbfc0-e642-4581-a96f-2e33ac17ee19, session:f7301d22-fd90-4675-95cd-17b006262ef9, artifact:6c552b3c-44e4-4b57-8b93-e5a645d057d1:5e15b337-adf1-483c-8ba6-7cd2de5f500b, artifact:7c70ca6e-3fb2-4cda-b75a-72702ee10d97:516c2288-3458-40b2-9e02-87a918c35a3a
\section*{Authorship note}
AI-assisted synthesis from saved research; not independently peer reviewed.
\end{document}
